\section{The detection pack and its telemetry}
\label{sec:pack}

\paragraph{Telemetry.}
The demo application emits one JSON object per event, aligned with the OpenTelemetry GenAI semantic
conventions~\cite{otel-genai}: \texttt{prompt}, \texttt{guardrail}, \texttt{retrieval},
\texttt{tool\_call}, \texttt{completion}, \texttt{memory\_write} and \texttt{tool\_registry}. Beyond
what an LLM gateway usually logs, it records four fields that most of the agentic rules depend on:
the provenance of every retrieved chunk (\texttt{doc\_source}), the normalised destination, risk class
and human-approval state of every tool call (\texttt{tool\_target}, \texttt{tool\_risk},
\texttt{tool\_approval}), a preview of every tool result (\texttt{tool\_result}, the first 1{,}000
characters), and, for every MCP tool the client discovers, the advertised description and its SHA-256
(\texttt{tool\_registry}). The MCP integration is real: the agent connects to a third-party server
over stdio with the official Python SDK (version 2.3).

For the frontier phase two event types were added. A \texttt{reasoning} event carries the model's
reasoning text for each model round when the provider returns one; for Claude models this is always a
provider-written summary, never the raw chain of thought, and is labelled as such. A
\texttt{refusal} event records a refusal by the provider: the API stopped generation with a refusal
stop reason (for Claude, \texttt{stop\_reason: refusal} with a category) or a content filter blocked
the request. A refusal in the model's own words is not an event; it is text in the answer or the
reasoning summary that a rule (\llmd{023}) has to recognise.

\paragraph{Rules.}
The pack has 25 detections, each implemented as a Microsoft Sentinel scheduled analytic rule (KQL);
23 are also implemented as Wazuh rules, and the two newest (\llmd{025} and \llmd{026}) are Sentinel
only so far (Table~\ref{tab:rules} in the appendix). They fall into four families that matter
for the results:
\begin{itemize}
  \item \emph{Content signatures} match text: instruction-override phrasing in prompts (\llmd{001}),
        in retrieved documents (\llmd{002}) and in tool results (\llmd{021}); obfuscation
        (\llmd{011}); system-prompt extraction requests (\llmd{012}); secrets in outputs
        (\llmd{004}); hidden instructions in MCP tool descriptions (\llmd{015}).
  \item \emph{Planted and structural signals} need no knowledge of the attack text: a canary token
        in the system prompt (\llmd{003}), a changed MCP description hash (\llmd{014}), auto-loading
        images to external hosts (\llmd{005}), dangerous commands (\llmd{017}), credential paths
        (\llmd{018}), writes to agent configuration (\llmd{020}).
  \item \emph{Behavioural and correlation rules} reason about what the agent did: egress to external
        destinations (\llmd{006}), egress after untrusted content entered the context (\llmd{007}),
        tool-call bursts (\llmd{008}), token spikes (\llmd{009}), repeated guardrail blocks
        (\llmd{010}), memory writes with standing instructions (\llmd{013}), destructive tools without
        approval (\llmd{016}), and a private read followed by egress in one turn (\llmd{019}), the
        ``lethal trifecta'' of Willison~\cite{trifecta}.
  \item \emph{Refusal, reasoning and answer signals}, added for this study: a refusal (provider or
        model) in a session that pulled in external content and took no high-risk action
        (\llmd{023}); a reasoning summary that states intent to follow injected content
        (\llmd{024}) or notes that an instruction was injected (\llmd{025}); and an answer that
        names an internal document and an outside email address without warning the user
        (\llmd{026}, Section~\ref{sec:relay}).
\end{itemize}

\paragraph{Rule groups used in the analysis.}
For the frontier analysis the rules are grouped by what they read. INPUT rules scan what enters the
model (\llmd{001}, 002, 010, 011, 012, 014, 015, 021). OUTPUT rules read what the agent produced or
did, split into \emph{output-content} rules (\llmd{003}, 004, 005, 018) and \emph{behavioural} rules
(\llmd{006}, 007, 008, 013, 016, 017, 019, 020). The \emph{refusal signal} is \llmd{023} or a
provider refusal event. The groups are defined in the pre-registration of the confirmatory run.

\paragraph{Prior validation.}
Before this study the first 21 rules were validated on a synthetic, labelled dataset of 70 scenarios:
all KQL rules match the labels exactly on the Kusto engine (Kustainer) and on a live Sentinel
workspace, and the Wazuh rules match except for two documented engine limits. That dataset was
written by the author of the rules; this paper replaces it with third-party data and live models.
